%++++++++++++++++++++++++++++++++++++++++
% Don't modify this section unless you know what you're doing!
\documentclass[letterpaper,11pt]{article}
\usepackage{natbib}
\bibliographystyle{unsrtnat}
\usepackage{tabularx} % extra features for tabular environment
\usepackage{amsmath}  % improve math presentation
\usepackage{graphicx} % takes care of graphic including machinery
\usepackage[margin=1in,letterpaper]{geometry} % decreases margins
% \usepackage{mathpazo}
\usepackage{libertine}
%\usepackage{cite} % takes care of citations
\usepackage[final]{hyperref} % adds hyper links inside the generated pdf file
\hypersetup{
	colorlinks=true,       % false: boxed links; true: colored links
	linkcolor=blue,        % color of internal links
	citecolor=blue,        % color of links to bibliography
	filecolor=magenta,     % color of file links
	urlcolor=blue         
}
%++++++++++++++++++++++++++++++++++++++++


\begin{document}

\title{Annual Progress Report 20  -20  }
\author{Name of the Student\\ HBNI Enrolment number: MATH123410000}
\date{\today}
\maketitle


\section{Progress made during academic year 20  -20  }
\subsection{Task 1}
Description

\subsection{Task 2}
Description

\begin{table}
    \centering
    \caption{Table}
    \begin{tabular}{|p{0.1\linewidth} | p{0.2\linewidth}|}
    \hline
     Category    &  Method  \\
     \hline
     One  &  Description  \\
     \hline
     Two & Description \\
     \hline
     Three & Description \\
     \hline
     Four  & Description \\
     \hline
    \end{tabular}
    \label{tab:methods}
\end{table}

\subsection{Task 3}
Description

\subsection{Task 4}
Description


\section{Papers published/communicated}

\begin{itemize}
    \item Published
    \begin{enumerate}
        \item \textbf{Title of the paper}          \href{google.com}{[preprint]}\\
        \textbf{author1}, author2, author3, and author4 \\
        \textit{Journal/Conference/Book, 2020}
        \item \textbf{Title of the paper}          \href{google.com}{[preprint]}\\
        \textbf{author1}, author2, author3, and author4 \\
        \textit{Journal/Conference/Book, 2020}
    \end{enumerate}
    \item Communicated
    \begin{enumerate}
        \item \textbf{Title of the paper}          \href{google.com}{[preprint]}\\
        \textbf{author1}, author2, author3, and author4 \\
        \textit{Journal/Conference/Book, 2020}
        \item \textbf{Title of the paper}          \href{google.com}{[preprint]}\\
        \textbf{author1}, author2, author3, and author4 \\
        \textit{Journal/Conference/Book, 2020}
    \end{enumerate}
\end{itemize}

\section{Courses}
\subsection{Courses: Completed}
\begin{enumerate}
    \item Course Name and Details
    \item Course Name and Details
\end{enumerate} 

\subsection{Courses: Teaching Assistant}
\begin{enumerate}
    \item \textbf{M123}: Course Name, Institute
    \item \textbf{M124}: Course Name, Institute
\end{enumerate}

\section{Work planned for next year}
\subsection{Goal 1}
Description

\subsection{Goal 2}
Description

\subsection{Goal 3}
Description

% \section{Introduction}

% Galaxies are very interesting. Measuring their properties is very important. 
% We even have an equation (see equation \ref{eq1})

% \begin{equation} \label{eq1} % the label is used to reference the equation
% V=\frac{8\pi a^{-5}}{3\lambda}
% \end{equation}

% where a is the distance to the Sun, $\lambda$ is something else.... Don't forget to
% explain what each variable means the first time that you introduce it.

% This is also where you explain the aims of the experiment.
% \section{Method}

% Describe what you did (not the names of files and directories of course, but all the important steps that bring you to the results). Why did you do this? did you have to make some choices along the way --- explain them. Did it not work well the first time, and you had to improve something? write about it

% Include diagrams/photos of the experimental setup (see Figure \ref{fig1}), or screenshots highlighting important steps of the process.

% \begin{figure}[!h] 
%         \centering \includegraphics[width=0.5\columnwidth]{Laser.jpg}
%         \caption{\label{fig1}Every figure MUST have a caption.
%         }
% \end{figure}

% \section{Results and Analysis}

% In this section you will need to show your results. Use tables and
% figures (long tables can be put in an Appendix). Table~\ref{table1} is an example.

% \begin{table}[ht]
% \begin{center}
% \caption{Every table needs a caption.}
% \label{table1} 
% \begin{tabular}{cc} 
% \hline
% \multicolumn{1}{c}{distance (m)} & \multicolumn{1}{c}{V (km s$^-1$)} \\
% \hline
% 0.0044151 &   0.0030871 \\
% 0.0021633 &   0.0021343 \\
% 0.0003600 &   0.0018642 \\
% 0.0023831 &   0.0013287 \\
% \hline
% \end{tabular}
% \end{center}
% \end{table}

% Here explain how you determine uncertainties for different
% measured values.

% If in the process of data analysis you found any noticeable systematic
% error(s), you have to explain them in this section of the report.


% \section{Discussion and Conclusion}
% Here you briefly summarize your findings. Do your results match published values? If not, what could the source(s) of errors? How could the experiment be improved?

\bibliography{ref}


% \appendix
% \section*{Appendix: Velocity measurements}
% You can put here long tables 
% \begin{table}[ht]
% \begin{center}
% \caption{Every table needs a caption.}
% \label{table2} 
% \begin{tabular}{cc} 
% \hline
% \multicolumn{1}{c}{distance (m)} & \multicolumn{1}{c}{V (km s$^-1$)} \\
% \hline
% 0.0044151 &   0.0030871 \\
% 0.0021633 &   0.0021343 \\
% 0.0003600 &   0.0018642 \\
% 0.0023831 &   0.0013287 \\
% 0.0044151 &   0.0030871 \\
% 0.0021633 &   0.0021343 \\
% 0.0003600 &   0.0018642 \\
% 0.0023831 &   0.0013287 \\
% 0.0044151 &   0.0030871 \\
% 0.0021633 &   0.0021343 \\
% 0.0003600 &   0.0018642 \\
% 0.0023831 &   0.0013287 \\
% 0.0044151 &   0.0030871 \\
% 0.0021633 &   0.0021343 \\
% 0.0003600 &   0.0018642 \\
% 0.0023831 &   0.0013287 \\
% 0.0044151 &   0.0030871 \\
% 0.0021633 &   0.0021343 \\
% 0.0003600 &   0.0018642 \\
% 0.0023831 &   0.0013287 \\
% 0.0044151 &   0.0030871 \\
% 0.0021633 &   0.0021343 \\
% 0.0003600 &   0.0018642 \\
% 0.0023831 &   0.0013287 \\
% \hline
% \end{tabular}
% \end{center}
% \end{table}

\end{document}
